\section{Marginal tails and an initial box estimate}
\label{sec:initial}

This section supplies two inputs to the proof: a moment bound for marginal
spectra and an entropy bound strictly below the volume of a rectangle that
stays away from the prescribed cut.  The spectral gap first controls the
fluctuations of the negative logarithms of marginal eigenvalues.  To bound
the entropy itself, we then optimize a product of positive operators on
nested rectangles.  An energy estimate and a trial using a Schmidt vector
give opposite bounds on the norm of the filtered ground vector.  Comparing
them discounts the conditional entropy of the innermost rectangle.
Averaging the entropy chain rule over
orders of small rectangles turns this discount into a subvolume bound.

\subsection{A marginal spectral estimate}

Boundary-dependent concentration of entanglement spectra was established
by Anshu, Harrow, and Soleimanifar through smooth entanglement spread,
with square-root boundary scaling on lattices up to logarithmic
factors~\cite{AnshuHarrowSoleimanifar2022}.  The estimate below instead
controls the moment-generating function of the marginal surprisal
centered at its entropy.  We prove it directly by tilting the ground
vector and comparing its excitation energy with the global gap.

The following lemma does not require a lattice.  Its dependence on the
logarithm of each interaction's supporting dimension will also allow us
to apply it to the truncated interactions constructed in
Section~\ref{sec:quasilocal}.  A related precedent is the bound on local
entanglement production whose constant is independent of spectator
dimensions~\cite{VanAcoleyenMarienVerstraete2013}.  Here the required
estimate is a quadratic bound on the real part of the change under
conjugation; we prove it below.

\begin{lemma}[Marginal moment and tail bounds]
\label{lem:tail}
Let $\mathcal H=\bigotimes_{v\in V}\mathcal H_v$ be a finite tensor product
of finite-dimensional Hilbert spaces.  Let
$\mathsf H=\sum_{i\in\mathcal I} h_i$, where $\mathcal I$ is finite,
each $h_i$ is self-adjoint, $\norm{h_i}\le c_0$ for a fixed $c_0\ge0$,
and $h_i$ has a designated support $D_i\subseteq V$.  Suppose that
$\mathsf H$ has a one-dimensional ground space spanned by a normalized
vector $\Psi$, ground energy $E_0$, and gap at least $g_0>0$.  Fix a
bipartition $B\sqcup B^c=V$.  Let $\mathcal I_\times$ be the indices
whose designated support meets both sides, and put
\[
 \rho=\Tr_{B^c}\ket{\Psi}\bra{\Psi},\qquad
 d_i=\prod_{v\in D_i}\dim\mathcal H_v,\qquad
 \mathcal B=1+\sum_{i\in\mathcal I_\times}\log^2(\mathrm e d_i).
\]
Let $p_j$ be the positive eigenvalues of $\rho$, let $K$ take the value
$-\log p_j$ with probability $p_j$, and write $S=S(\rho)$ and
$\vartheta=c_0/g_0$.  Then
\begin{equation}
 \log\E\mathrm e^{uK}
 \le uS+512\mathrm e\,\vartheta\mathcal B u^2,
 \qquad
 |u|\le\frac{1}{32\sqrt{(1+\vartheta)\mathcal B}}.
 \label{eq:initial-tail-mgf}
\end{equation}
Consequently, for every $w\ge0$,
\begin{equation}
 \Pr\{|K-S|>w\}
 \le \min\left\{1,\,
  2\mathrm e^{\mathrm e/2}
  \exp\left(-\frac{w}{32\sqrt{(1+\vartheta)\mathcal B}}\right)
              \right\}.
 \label{eq:initial-tail-probability}
\end{equation}
\end{lemma}

\begin{proof}
Put $M(u)=\sum_jp_j^{1-u}$ and $L(u)=\log M(u)$.  Let $P$ be the
support projection of $\rho$.  On $P\mathcal H_B$ define
$F=\rho^{-u/2}$, and extend $F$ by the identity on the kernel of $\rho$.
The normalized tilted vector and its marginal are
\[
 \Theta=\frac{F\Psi}{\sqrt{M(u)}},\qquad
 \sigma=\Tr_{B^c}\ket{\Theta}\bra{\Theta}
       =\frac{\rho^{1-u}}{M(u)},
\]
where the last power is zero on the kernel.  Since $F$ is invertible on
the whole space,
\begin{equation}
 F\mathsf H F^{-1}\Theta=E_0\Theta.
 \label{eq:initial-exact-tilt}
\end{equation}
An interaction supported in $B^c$ is unchanged by this conjugation.  An
interaction supported in $B$ has unchanged expectation, because
\[
 \Tr\sigma Fh_iF^{-1}=\Tr\sigma h_i
\]
by $[F,\sigma]=0$ and cyclicity.  Thus only crossing interactions
contribute to the excitation energy of $\Theta$.

Fix such an interaction.  For every complex $z$, use the support
convention
\[
 \sigma^z=P\exp\bigl(z\log(\sigma|_{P\mathcal H_B})\bigr)P,
 \qquad
 f_i(z)=\langle\Theta,\sigma^z h_i\sigma^{-z}\Theta\rangle.
\]
This is an entire scalar function.  Although $\sigma^0=P$, the vector
$\Theta$ lies in its support, so $f_i(0)=\langle\Theta,h_i\Theta\rangle$.
At
\[
 z_u=-\frac{u}{2(1-u)}
\]
the two scalar normalization factors cancel, giving
\begin{equation}
 f_i(z_u)=\langle\Theta,Fh_iF^{-1}\Theta\rangle.
 \label{eq:initial-support-conjugation}
\end{equation}
Indeed, on the support,
$\sigma^{z_u}=M(u)^{-z_u}F$ and
$\sigma^{-z_u}=M(u)^{z_u}F^{-1}$.  The endpoint support projections
in the expectation remove any matrix elements entering the kernel.

On the imaginary axis the powers are unitary on the support, and hence
$|f_i(\mathrm i y)|\le c_0$.  Decompose $h_i$ across the cut by matrix
units on its portion in $B$:
\[
 h_i=\sum_{\alpha=1}^{N_i}c_\alpha\otimes d_\alpha,
 \qquad N_i\le d_i^2,
 \qquad \norm{c_\alpha}\norm{d_\alpha}\le c_0.
\]
Here each factor acts as the identity outside the appropriate portion
of $D_i$.  If $s_j$ are the positive Schmidt probabilities of $\Theta$,
a product contributes terms with scalar coefficient
$s_j^{1/2+z}s_k^{1/2-z}$.  For $|\Re z|\le1/4$,
\[
 \bigl|s_j^{1/2+z}s_k^{1/2-z}\bigr|\le s_j+s_k.
\]
Moreover, in the two Schmidt bases, Cauchy--Schwarz bounds every row
and column sum of $|(c_\alpha)_{jk}(d_\alpha)_{jk}|$ by
$\norm{c_\alpha}\norm{d_\alpha}$.  Compression to the Schmidt supports
does not increase either norm.  Summing first over a row or a column
therefore gives
\[
 |f_i(z)|\le2c_0d_i^2\qquad(|\Re z|\le1/4).
\]
These bounds are uniform in $\Im z$.  Apply the three-lines theorem
on each of the two strips between real parts $0$ and $\pm1/4$.
Writing $\ell_i=\log(\mathrm e d_i)$, we obtain
\begin{equation}
 |f_i(z)|\le\mathrm e c_0
 \qquad\left(|\Re z|\le\delta_i:=\frac1{8\ell_i}\right),
 \label{eq:initial-narrow-strip}
\end{equation}
because $\log(2d_i^2)\le2\ell_i$.  The case $c_0=0$ is immediate;
otherwise this application may be made to $f_i/c_0$.

Hermiticity gives
$f_i(-\overline z)=\overline{f_i(z)}$.  The even Taylor coefficients
at zero are therefore real, and the odd coefficients are purely
imaginary.  Cauchy's coefficient estimate on the disk of radius
$\delta_i$, followed by summation of the even terms, yields, for real
$|x|\le\delta_i/2$,
\begin{equation}
 \bigl|\Re(f_i(x)-f_i(0))\bigr|
 \le2\mathrm e c_0x^2\delta_i^{-2}
 =128\mathrm e c_0\ell_i^2x^2.
 \label{eq:initial-real-quadratic}
\end{equation}
In particular, it is the real linear term that vanishes; we do not
assert that $f_i'(0)=0$.

For $u$ in the interval of the lemma, $|u|\le1/32$ and
$|z_u|\le|u|$.  Since $\ell_i\le\sqrt{\mathcal B}$, the estimate
\eqref{eq:initial-real-quadratic} applies to every crossing interaction.
Using \eqref{eq:initial-exact-tilt} and
\eqref{eq:initial-support-conjugation}, we conclude that
\begin{equation}
 0\le\langle\Theta,(\mathsf H-E_0)\Theta\rangle
 \le128\mathrm e c_0\mathcal B u^2.
 \label{eq:initial-tilted-energy}
\end{equation}
The full-system gap inequality
$\mathsf H-E_0\ge g_0(\id-\ket{\Psi}\bra{\Psi})$ now gives
\[
 \frac{M(u/2)^2}{M(u)}
 =|\langle\Psi,\Theta\rangle|^2
 \ge1-128\mathrm e\,\vartheta\mathcal B u^2.
\]
The last subtracted quantity is at most
$\mathrm e\vartheta/[8(1+\vartheta)]<1/2$.  Taking logarithms and
using $-\log(1-v)\le2v$ for $0\le v\le1/2$ gives
\[
 L(u)\le2L(u/2)+256\mathrm e\,\vartheta\mathcal B u^2.
\]
Iterate this inequality.  Every successive halving remains in the
permitted interval, for either sign of $u$, and
$2^jL(u/2^j)\to uS$ because $L(0)=0$ and $L'(0)=S$.
The sum of the remaining geometric series is at most two, proving
\eqref{eq:initial-tail-mgf}.  Finally apply the exponential Markov
inequality with
$u=\pm[32\sqrt{(1+\vartheta)\mathcal B}]^{-1}$.
The quadratic term in \eqref{eq:initial-tail-mgf} is at most
$\mathrm e/2$.  Adding the two one-sided estimates proves
\eqref{eq:initial-tail-probability}.
\end{proof}

For the finite-range lattice Hamiltonians of this paper, the interaction
budget and support-size bounds of Section~\ref{sec:prelim} imply
\begin{equation}
 \mathcal B\le C(1+|\partial_\Lambda B|).
 \label{eq:initial-local-tail-parameter}
\end{equation}
Thus, for every $n\ge1$ and fixed $M>0$, the negative logarithms of the
positive eigenvalues of a marginal at a cut with $b'$ edges lie within
$C_M\sqrt{1+b'}\log(n+2)$ of its entropy except for probability
at most $(n+2)^{-M}$.  All constants in this
specialization depend only on $q,R,J,\Delta$ and, for $C_M$, on $M$.

This concentration estimate controls fluctuations around the entropy,
but does not yet bound that center.  The next argument compares two
estimates for the same norm to obtain a conditional-entropy discount,
from which the rectangle bound will follow.

\subsection{An entropy discount from a buffer}

Fix the set $A\subseteq\Lambda$ whose entropy will ultimately be
bounded, and let $Z$ be the set of endpoints of its crossing edges.
For an integer rectangle $Q$---a product of two finite nonempty integer
intervals---its size is the larger number of sites along an axis.
Fix $D_0>2R+10$.  We call $Q$ \emph{safe} if
\[
 \dist_\infty(Q,Z)>D_0\,\operatorname{size}(Q),
\]
with distance to the empty set taken to be infinite.  For an integer
$d\ge0$ write
\[
 Q^{+d}=Q+\bigl([-d,d]^2\cap\Z^2\bigr).
\]
These are ambient rectangles; all physical regions below are their
intersections with $A$.  In particular, their use does not change the
graph-distance locality of the Hamiltonian.  For a unit vector $\xi$,
we write $\rho_{\xi,D}$ for its marginal on $D$ and $S_\xi(D)$ for
that marginal's entropy.

\begin{lemma}[Conditional entropy in a buffered rectangle]
\label{lem:initial-buffer}
There are integers $C_{\mathrm{pad}}\ge1$ and a constant
$C_{\mathrm{buf}}$, depending only on $q,R,J,\Delta$, with the following
property.  Let $Q$ be safe, let $r\ge1$ be an integer, and let $Q_0$ be
an integer rectangle of size at most $r$ such that
$Q_0^{+C_{\mathrm{pad}}r}\subseteq Q$.  Set
\[
 X=A\cap Q_0,\qquad
 T=(A\cap Q_0^{+C_{\mathrm{pad}}r})\setminus X.
\]
Then, in the original ground state $\Omega$,
\begin{equation}
 S(X\mid T)\le\frac12S(X)+C_{\mathrm{buf}}r.
 \label{eq:initial-buffer-discount}
\end{equation}
\end{lemma}

\begin{proof}
If $X$ is empty, the conclusion is immediate.  Suppose henceforth that
$X\ne\varnothing$.  We will compare two estimates for the largest norm
obtained by applying positive filters on nested regions around $X$.
The filters need not commute.  Stationarity and the gap give a comparison
involving the entropies of those regions.  For the opposite comparison,
we pin a Schmidt vector of $X$ and then average over its index.  This
removes $X$ from each region and contributes a single entropy cost $S(X)$.
The difference of the two comparisons will therefore control a weighted
sum of conditional entropies of $X$.

Choose an integer $D>2R+2$.  Temporarily leave the integer
$C_{\mathrm{pad}}>D+1$ variable, and put
\[
 m=\left\lfloor\frac{(C_{\mathrm{pad}}-1)r}{D}\right\rfloor,
 \qquad d_j=r+jD,\qquad
 X_j=A\cap Q_0^{+d_j}\quad(1\le j\le m),\qquad X_0=X.
\]
All these rectangles lie in the safe parent $Q$.  Any interaction
meeting an $X_j$ lies wholly in $A$: otherwise a graph path of length
at most $R$ from a support site in $X_j$ to a support site of the
opposite color would contain a point of $Z$ within ambient distance
$R$ of $Q$.  This contradicts safety.  The same observation rules out
any nearest-neighbor color-crossing edge incident with $X_j$.
Consequently the edge boundary of $X_j$ has at most
$4(r+2d_j)$ edges, all crossing the ambient rectangle's perimeter.
Missing sites in $\Lambda$ only reduce this count.

Every designated interaction support splits at most one of the cuts
$X_j\mid X_j^c$.  More precisely, a support split by $X_j$ is disjoint
from every $X_i$ with $i<j$ and contained in every $X_i$ with $i>j$.
Indeed, successive rectangles have an ambient margin $D>R$, whereas
the ambient diameter of the support is at most its graph diameter
$R$.  A term that does not split a contour is either contained in it
or disjoint from it.  The support budget and
\eqref{eq:initial-local-tail-parameter} therefore give, with constants
independent of $C_{\mathrm{pad}}$,
\begin{equation}
 \#\{\text{interactions split by }X_j\}\le C(r+d_j),
 \qquad \mathcal B_j\le C(r+d_j),
 \label{eq:initial-contour-budget}
\end{equation}
where $\mathcal B_j$ is the parameter in Lemma~\ref{lem:tail} for the
marginal on $X_j$.

Assign a positive weight $a_j$ to each contour.  The energy estimate below
charges contour $j$ by a constant times
$(r+d_j)a_j^2$.  We choose weights of total mass two: in the eventual
conditional-entropy comparison this will give the factor two needed for
\eqref{eq:initial-buffer-discount}.  At this fixed mass, Cauchy--Schwarz
gives
\[
 4=\left(\sum_j a_j\right)^2
 \le\left(\sum_j(r+d_j)a_j^2\right)
       \left(\sum_j\frac1{r+d_j}\right).
\]
Equality holds for weights proportional to $1/(r+d_j)$.  Accordingly, put
\begin{equation}
 Z_r=\sum_{j=1}^m\frac1{r+d_j},\qquad
 a_j=\frac{2}{Z_r(r+d_j)},\qquad
 \mathcal E=\sum_{j=1}^m(r+d_j)a_j^2=\frac4{Z_r}.
 \label{eq:initial-harmonic-weights}
\end{equation}
Their sum is two, and their quadratic cost attains this lower bound.
Integral comparison gives the uniform lower bound
\begin{equation}
 Z_r\ge\frac1D\log\frac{C_{\mathrm{pad}}+1}{2+D}.
 \label{eq:initial-harmonic-lower}
\end{equation}
For example, compare the sum with the integral of
$(2r+Dx)^{-1}$ over $1\le x\le m+1$, use
$2r+D(m+1)>(C_{\mathrm{pad}}+1)r$, and use $r\ge1$ in the denominator.
Thus we may make $\mathcal E$ as small as needed, uniformly in $r$.
We take $C_{\mathrm{pad}}$ large enough that $a_j\le1/2$ for all $j$.
We shall increase it once more after the energy estimate below; all
constants up to that choice are independent of $C_{\mathrm{pad}}$.

\paragraph{Stationarity of the optimized filters.}
Fix a sufficiently small $\epsilon>0$, and maximize
\[
 N=\norm{L_m\cdots L_1\Omega}
\]
over positive operators $L_j$ on $X_j$ satisfying
\begin{equation}
 L_j\ge\epsilon\id,
 \qquad \Tr L_j^{2/a_j}=1.
 \label{eq:initial-filter-constraints}
\end{equation}
The trace is over $\mathcal H_{X_j}$.  Choose $\epsilon$ so small that
$(\dim\mathcal H_{X_j})\epsilon^{2/a_j}<1$ for every $j$.
The feasible sets are nonempty and compact.  A maximum exists and is
positive, since every feasible product is invertible.  At an optimum
write
\[
 K=L_m\cdots L_1,\qquad \phi=K\Omega/N.
\]
We claim that
\begin{equation}
 [L_j,\rho_{\phi,X_j}]=0\qquad(1\le j\le m).
 \label{eq:initial-filter-commutation}
\end{equation}
Prove this in descending order.  For an anti-Hermitian operator $B$ on
$X_j$, vary $L_j$ to $\mathrm e^{tB}L_j\mathrm e^{-tB}$.
This preserves both constraints in
\eqref{eq:initial-filter-constraints}.  With
$K_{>j}=L_m\cdots L_{j+1}$, the derivative of the logarithmic norm is
\[
 \Re\langle\phi,
 K_{>j}(B-L_jBL_j^{-1})K_{>j}^{-1}\phi\rangle.
\]
Assume that \eqref{eq:initial-filter-commutation} is already known for
indices greater than $j$.  In this expectation the outer factor $L_m$
can be removed by cyclicity of the trace on $X_m$, because its marginal
commutes with $L_m$ and the operator inside is supported on $X_m$.
Repeat this argument down to $L_{j+1}$.  Stationarity, together with
the fact that the expectation of $B$ is purely imaginary, gives
\[
 \Re\Tr(L_j^{-1}\rho_{\phi,X_j}L_jB)=0
 \quad\text{for every anti-Hermitian }B.
\]
The matrix $L_j^{-1}\rho_{\phi,X_j}L_j$ is therefore Hermitian.
The equality to its adjoint, multiplied on both sides by $L_j$,
shows that $\rho_{\phi,X_j}$ commutes with $L_j^2$, hence with
$L_j$.  This proves the induction.  Only expectations have been
simplified; no commutation between distinct filters is asserted.

In a common eigenbasis let $l_i$ be the eigenvalues of $L_j$, let $p_i$
be those of $\rho_{\phi,X_j}$, and set $x_i=l_i^{2/a_j}$.
A diagonal variation and the same descending trace argument give
\[
 d\log N=\frac{a_j}{2}\sum_i p_i\frac{dx_i}{x_i},
 \qquad \sum_i x_i=1,\qquad x_i\ge\epsilon^{2/a_j}.
\]
The ratios $p_i/x_i$ are equal to a common value $\lambda$ on free
coordinates and are at most $\lambda$ on coordinates at the floor.
There is a free coordinate by strict dimensional slack.  Moreover
$\lambda>0$: if all free coordinates had zero probability, transfer
from one of them to a positive-probability floored coordinate would
increase the norm.  Consequently
\begin{equation}
 x_i=\max\{\epsilon^{2/a_j},p_i/\lambda\},\qquad
 \left|\log\frac{l_i}{l_k}\right|
 \le\frac{a_j}{2}\left|\log\frac{p_i}{p_k}\right|
 \quad(p_i,p_k>0).
 \label{eq:initial-filter-clipping}
\end{equation}
The last inequality follows because taking a maximum with a common
positive floor contracts differences of logarithms.

\paragraph{The optimized vector has small excitation energy.}
We have the exact identity $KHK^{-1}\phi=E_0\phi$ for the original
full Hamiltonian.  Consider one interaction in its expectation.
All outer filters on regions containing its support can be removed
by \eqref{eq:initial-filter-commutation} and cyclicity, as above.
All remaining inner filters whose regions are disjoint from its
support commute with the interaction.  The single-split property
leaves at most one conjugation $L_jh_iL_j^{-1}$.

In Schmidt bases for $X_j\mid X_j^c$, Hermitian pairing replaces the
real part of this conjugation by coefficients involving
\[
 \sqrt{p_i p_k}
 \left(\cosh\left(\log\frac{l_i}{l_k}\right)-1\right).
\]
For positive $p_i,p_k$, \eqref{eq:initial-filter-clipping} and
$a_j\le1/2$ bound this expression by $C a_j^2(p_i+p_k)$.
To verify the uniform constant, put
$z=\tfrac12\log(p_i/p_k)$ and use
\[
 \cosh(a_jz)-1\le\tfrac12a_j^2z^2\cosh(a_jz),
 \qquad
 \sup_{z\in\R}\frac{z^2\cosh(z/2)}{\cosh z}<\infty.
\]
Pairs containing a zero probability have zero coefficient.
A finite-range interaction has a product decomposition across this
cut with a uniformly bounded sum of norm products.  The row and
column Cauchy--Schwarz estimate used in Lemma~\ref{lem:tail} now bounds
its contribution by $C a_j^2$.  Summing
\eqref{eq:initial-contour-budget} gives
\begin{equation}
 0\le\langle\phi,(H-E_0)\phi\rangle\le C\mathcal E.
 \label{eq:initial-product-energy}
\end{equation}
The constant depends only on $q,R,J$, not on the number of filters or
their supports' dimensions.

The full-system gap converts \eqref{eq:initial-product-energy} into
$1-|\langle\Omega,\phi\rangle|^2\le C\mathcal E/\Delta$.
Choose $C_{\mathrm{pad}}$ large enough that this is at most $1/4$.
At the same time arrange that
\begin{equation}
 \frac{a_j}{1-a_j}
 \le\frac{1}{32\sqrt{(1+J/\Delta)\mathcal B_j}}
 \quad\text{for every }j.
 \label{eq:initial-weight-range}
\end{equation}
Both choices are uniform in $r$: by
\eqref{eq:initial-contour-budget}, the right side is bounded below
by a parameter-dependent constant times $(r+d_j)^{-1/2}$, whereas
the left side is at most $4/[Z_r(r+d_j)]$.
Equation~\eqref{eq:initial-harmonic-lower} completes the choice of
$C_{\mathrm{pad}}$, which is now fixed.

\paragraph{The first norm comparison.}
Let $N_j$ be the optimum for the prefix of filters $1,\ldots,j$ at the
same floor, and set $N_0=1$.  The preceding stationarity and energy
arguments apply separately to each prefix: its summed cost is at
most $\mathcal E$.  Thus, for a parameter-dependent constant $C$,
its normalized output has overlap at least
$\sqrt{1-C\mathcal E}$ with $\Omega$, where
$C\mathcal E\le1/4$.  For an optimizing prefix, put $K_{<j}=L_{j-1}\cdots L_1$.
Cauchy--Schwarz gives
\begin{align*}
 \sqrt{1-C\mathcal E}\,N_j
 &\le |\langle\Omega,L_jK_{<j}\Omega\rangle|
   =|\langle L_j\Omega,K_{<j}\Omega\rangle|\\
 &\le\norm{L_j\Omega}\,N_{j-1}.
\end{align*}
For any feasible $L_j$, Schatten H\"older gives
\[
 \norm{L_j\Omega}^2
 =\Tr\rho_{\Omega,X_j}L_j^2
 \le\left(\Tr\rho_{\Omega,X_j}^{1/(1-a_j)}\right)^{1-a_j}.
\]
Apply Lemma~\ref{lem:tail} with
$u=-a_j/(1-a_j)$, permitted by
\eqref{eq:initial-weight-range}.  We obtain
\[
 \log\norm{L_j\Omega}^2
 \le-a_jS(X_j)+C(r+d_j)a_j^2.
\]
After summing logarithms over the prefixes and using
$\log(1-C\mathcal E)\ge-C'\mathcal E$, this gives
\begin{equation}
 -2\log N\ge\sum_{j=1}^m a_jS(X_j)-C(m+1)\mathcal E.
 \label{eq:initial-norm-lower}
\end{equation}
The constants in this estimate are independent of the floor.

\paragraph{The opposite comparison by Schmidt pinning.}
Let the floor decrease to zero at this fixed finite instance.
The optima converge to the optimum over nonnegative filters with the
same trace constraints.  Indeed, these latter feasible sets are
compact, and every feasible nonnegative operator is a limit of
strictly positive, trace-normalized operators.  Taking their minimal
eigenvalues as admissible floors proves the assertion.
We continue to denote the limiting optimum by $N$;
\eqref{eq:initial-norm-lower} remains valid.

Write the Schmidt decomposition at $X$ as
\[
 \Omega=\sum_i\sqrt{p_i}\,u_i\otimes w_i,
 \qquad p_i>0,
\]
and let $T_j=X_j\setminus X$.  For a fixed index $i$ use trial filters
\[
 L_j=\ket{u_i}\bra{u_i}\otimes M_j,
 \qquad M_j\ge0\text{ on }T_j,
 \qquad\Tr M_j^{2/a_j}=1.
\]
They are feasible in the zero-floor limit and give
$N\ge\sqrt{p_i}\,N_i$, where
\[
 N_i=\sup\norm{M_m\cdots M_1w_i}
\]
under the displayed constraints.  We claim that
\begin{equation}
 \log N_i\ge-\frac12\sum_j a_jS_{w_i}(T_j).
 \label{eq:initial-trial-norm}
\end{equation}
For fixed invertible choices $M_j$, consider
$F_i(z)=M_m^z\cdots M_1^z w_i$ on $0\le\Re z\le1$.
It is bounded and analytic there.  Its norm on the imaginary axis is
one.  On the other boundary, put $U_j=M_j^{\mathrm i y}$ and
$V_j=U_j\cdots U_1$, with $V_0=\id$.  Moving these unitaries to the
left gives the exact identity
\[
 M_m^{1+\mathrm i y}\cdots M_1^{1+\mathrm i y}
 =(U_m\cdots U_1)
   (V_{m-1}^*M_mV_{m-1})\cdots(V_0^*M_1V_0).
\]
Each positive factor is conjugated only by unitaries supported on
smaller nested regions.  The remaining factors are
therefore still positive, supported on their respective $T_j$, and
have the same trace-power normalizations.  The leftmost product of
unitaries does not change the norm, so
$\norm{F_i(1+\mathrm i y)}\le N_i$.
The vector form of the three-lines theorem, obtained by testing
against unit vectors and applying the scalar lemma
\cite[Lemma~5.10]{Tao247ANotes1}, gives
$\log\norm{F_i(x)}\le x\log N_i$ for $0\le x\le1$.
Differentiating from the right at zero yields
\[
 \sum_j\langle w_i,\log M_j\,w_i\rangle\le\log N_i.
\]
For $0<\delta<1$, take
\[
 M_j=\left((1-\delta)\rho_{w_i,T_j}
       +\frac{\delta}{\dim\mathcal H_{T_j}}\id\right)^{a_j/2}.
\]
These are strictly positive and satisfy $\Tr M_j^{2/a_j}=1$.
As $\delta\downarrow0$, their expectations of $\log M_j$ tend to
$-a_jS_{w_i}(T_j)/2$, since the marginal gives zero weight to its kernel.
This proves
\eqref{eq:initial-trial-norm}.

For each Schmidt index the trial norm now gives
\[
 -2\log N\le-\log p_i+\sum_j a_jS_{w_i}(T_j).
\]
Average with probabilities $p_i$ and use concavity of each $T_j$
entropy.  Their averaged states are the original marginals, so
\begin{equation}
 -2\log N\le S(X)+\sum_j a_jS(T_j).
 \label{eq:initial-norm-upper}
\end{equation}
The two norm estimates have now served their purpose.  Combining
\eqref{eq:initial-norm-lower} and \eqref{eq:initial-norm-upper} gives
\[
 \sum_j a_j S(X\mid T_j)\le S(X)+C(m+1)\mathcal E.
\]
Since the $T_j$ are nested and $\sum_j a_j=2$, strong subadditivity
bounds the left side below by $2S(X\mid T_m)$.
Finally $T_m\subseteq T$, so conditioning further decreases the
conditional entropy.  Because $m\le C_{\mathrm{pad}}r/D$ and all
parameters have now been fixed, the resulting error is at most
$C_{\mathrm{buf}}r$.  This proves
\eqref{eq:initial-buffer-discount}.
\end{proof}

\subsection{The first subvolume exponent}

For an integer $r\ge1$, let $F_{\mathrm{box}}(r)$ be the supremum of
$S_\Omega(A\cap Q)$ over all admissible Hamiltonians and cuts with
fixed $q,R,J,\Delta$, and all safe integer rectangles $Q$ of size at
most $r$.  This is a finite, nondecreasing function, since
$F_{\mathrm{box}}(r)\le r^2\log q$.

\begin{proposition}[Initial safe-box estimate]
\label{prop:initial-box}
There are constants $C<\infty$ and $e_0\in(0,1)$, depending only on
$q,R,J,\Delta$, such that
\begin{equation}
 F_{\mathrm{box}}(r)\le Cr^{1+e_0}\qquad(r\ge1).
 \label{eq:initial-box-bound}
\end{equation}
\end{proposition}

\begin{proof}
Choose an integer $M$ below, and consider a safe parent rectangle of
size at most $Mr$.  If its size is at most $r$, its entropy is already
bounded by $F_{\mathrm{box}}(r)$.  Otherwise partition each axis into
chunks of $r$ sites and, if needed, one final shorter chunk.  Their
Cartesian products give at most $M^2$ children $Q_i$, each safe and
of size at most $r$.  Put $X_i=A\cap Q_i$.

Call a child interior if its $C_{\mathrm{pad}}r$ padding lies in the
parent.  There are at most
\[
 B_M=4(C_{\mathrm{pad}}+2)M
\]
noninterior children: in either coordinate only the first or last
$C_{\mathrm{pad}}+2$ chunk positions can violate this containment.
This estimate remains valid when the parent is thin and all of its
children are noninterior.

For an interior child its padded rectangle meets at most
$(2C_{\mathrm{pad}}+5)^2-1$ other children.  Order all children by a
uniform random permutation.  The probability that this child comes
after every one of those other children is at least
\[
 p_0=(2C_{\mathrm{pad}}+5)^{-2}>0.
\]
On that event, its preceding physical sites contain the entire
buffer in Lemma~\ref{lem:initial-buffer}.  The corresponding term of
the entropy chain rule is therefore at most
$S(X_i)/2+C_{\mathrm{buf}}r$.  In all other cases it is at most
$S(X_i)$.  The latter bound also holds for noninterior children.
Taking the expected entropy chain rule consequently gives
\begin{align}
 F_{\mathrm{box}}(Mr)
 &\le\left[(1-p_0/2)M^2+(p_0/2)B_M\right]
          F_{\mathrm{box}}(r)+C M^2r.
 \label{eq:initial-box-recursion}
\end{align}
For completeness, the probability of the buffer event may vary among
children.  Its lower bound $p_0$ controls the negative term
$-p_0S(X_i)/2$, since $S(X_i)\ge0$; its positive error can always be
bounded by $C_{\mathrm{buf}}r$.  Thus the displayed common bound
holds without assuming equal event probabilities.  Parents of size
at most $r$ satisfy the same inequality once its coefficient is at
least one.

Choose $M\ge16(C_{\mathrm{pad}}+2)$ and $M\ge2$.  Then the coefficient
in \eqref{eq:initial-box-recursion} is at most
\[
 \lambda=(1-p_0/4)M^2,
 \qquad M<\lambda<M^2.
\]
Iterate the recursion at $r=1,M,M^2,\ldots$.  The inhomogeneous terms
are bounded by a geometric series with ratio $M/\lambda<1$, yielding
$F_{\mathrm{box}}(M^j)\le C\lambda^j$ for all $j\ge0$.
Set
\[
 1+e_0=\log_M\lambda
       =2+\frac{\log(1-p_0/4)}{\log M}\in(1,2).
\]
For a general $r\ge1$, comparison with the next power of $M$ and
monotonicity of $F_{\mathrm{box}}$ prove
\eqref{eq:initial-box-bound} after a fixed change of constant.
\end{proof}

The estimate just proved is uniform over induced domains, including
those with holes or disconnected components.  Its only energy
comparison was with the original full Hamiltonian.  In particular,
we have not required a gap for any Hamiltonian obtained by restricting
to a rectangle.  The following sections will improve the exponent
$e_0$ and then replace the box estimate by a bound on mutual information
across a collar.
